
\begin{document}

% @see : https://coopmaths.fr/alea?uuid=c8f4b&id=4C10-5&n=1&d=10&s=false&alea=i6Nq&cd=1&cols=1
\begin{EXO}{Calculer.}{4C10-5}

 $\def\arraystretch{1.5}\begin{array}{|c|c|c|c|c|}
    \hline
    \times & -7 & +3 & -6 & +9 \\
    \hline
    -9 &  &  & &  \\
    \hline
    +5 & & & & \\
    \hline
    -2 & & & & \\
    \hline
    +3 & & & & \\
    \hline
    \end{array}$
\end{EXO}

% @see : https://coopmaths.fr/alea?uuid=c8f4b&id=4C10-5&n=1&d=10&s=false&alea=cQ0A&cd=1&cols=1
\begin{EXO}{Calculer.}{4C10-5}

 $\def\arraystretch{1.5}\begin{array}{|c|c|c|c|c|}
    \hline
    \times & +7 & -4 & +8 & -3 \\
    \hline
    +3 &  &  & &  \\
    \hline
    -4 & & & & \\
    \hline
    +6 & & & & \\
    \hline
    -2 & & & & \\
    \hline
    \end{array}$
\end{EXO}

% @see : https://coopmaths.fr/alea?uuid=73187&id=4C10-6&n=4&d=10&s=1&alea=y1Vg&cd=0&cols=2
\begin{EXO}{}{4C10-6}
\begin{multicols}{2}
\begin{enumerate}[itemsep=2em]
	\item \begin{minipage}[t]{\linewidth} Donner le signe de $y$ pour que $A$ soit négatif. \\$ A =  y \times (+18)\times (+19)\times (+19) $ \\ 
    
    	$\square\;$ négatif\qquad $\square\;$ positif\qquad \\
 \end{minipage}
	\item \begin{minipage}[t]{\linewidth} Donner le signe de $x$ pour que $B$ soit positif. \\$ B =  x \times (-20)\times (-20)\times (+10) $ \\ 
    
    	$\square\;$ négatif\qquad $\square\;$ positif\qquad \\
 \end{minipage}
	\item \begin{minipage}[t]{\linewidth} Donner le signe de $a$ pour que $C$ soit négatif. \\$ C =  (+12) \times a\times (-12)\times (+4)\times (-20) $ \\ 
    
    	$\square\;$ négatif\qquad $\square\;$ positif\qquad \\
 \end{minipage}
	\item \begin{minipage}[t]{\linewidth} Donner le signe de $n$ pour que $D$ soit positif. \\$ D =  n \times (-9)\times (+9) $ \\ 
    
    	$\square\;$ négatif\qquad $\square\;$ positif\qquad \\
 \end{minipage}
\end{enumerate}
\end{multicols}
\end{EXO}

% @see : https://coopmaths.fr/alea?uuid=73187&id=4C10-6&n=4&d=10&s=2&alea=FLRe&cd=0&cols=2
\begin{EXO}{}{4C10-6}
\begin{multicols}{2}
\begin{enumerate}[itemsep=2em]
	\item \begin{minipage}[t]{\linewidth} Donner le signe de $n$ pour que $A$ soit négatif. \\$ A =  \dfrac {(+3)\times (-16)\times (+18)}{(+6)\times n\times (+13)} $ \\ 
    
    	$\square\;$ négatif\qquad $\square\;$ positif\qquad \\
 \end{minipage}
	\item \begin{minipage}[t]{\linewidth} Donner le signe de $m$ pour que $B$ soit positif. \\$ B =  \dfrac {(+14)\times m\times (-17)\times (-1)}{(+11)\times (-3)} $ \\ 
    
    	$\square\;$ négatif\qquad $\square\;$ positif\qquad \\
 \end{minipage}
	\item \begin{minipage}[t]{\linewidth} Donner le signe de $a$ pour que $C$ soit positif. \\$ C =  \dfrac {(+9)\times (+9)\times a}{(-17)\times (+2)\times (-17)} $ \\ 
    
    	$\square\;$ négatif\qquad $\square\;$ positif\qquad \\
 \end{minipage}
	\item \begin{minipage}[t]{\linewidth} Donner le signe de $m$ pour que $D$ soit positif. \\$ D =  \dfrac {(-11)\times (+18)\times m}{(+11)} $ \\ 
    
    	$\square\;$ négatif\qquad $\square\;$ positif\qquad \\
 \end{minipage}
\end{enumerate}
\end{multicols}
\end{EXO}

% @see : https://coopmaths.fr/alea?uuid=73187&id=4C10-6&n=4&d=10&s=3&alea=Vbf7&cd=0&cols=2
\begin{EXO}{}{4C10-6}
\begin{multicols}{2}
\begin{enumerate}[itemsep=2em]
	\item \begin{minipage}[t]{\linewidth} Donner le signe de $a$ pour que $A$ soit positif. \\$ A =  \dfrac {(-5)\times a\times (-17)}{(+4)} $ \\ 
    
    	$\square\;$ négatif\qquad $\square\;$ positif\qquad \\
 \end{minipage}
	\item \begin{minipage}[t]{\linewidth} Donner le signe de $m$ pour que $B$ soit positif. \\$ B =  (-5) \times (+12)\times (+15)\times m\times (+6) $ \\ 
    
    	$\square\;$ négatif\qquad $\square\;$ positif\qquad \\
 \end{minipage}
	\item \begin{minipage}[t]{\linewidth} Donner le signe de $n$ pour que $C$ soit négatif. \\$ C =  (+4) \times n\times (+20)\times (+7) $ \\ 
    
    	$\square\;$ négatif\qquad $\square\;$ positif\qquad \\
 \end{minipage}
	\item \begin{minipage}[t]{\linewidth} Donner le signe de $n$ pour que $D$ soit négatif. \\$ D =  \dfrac {(-1)\times n\times (-17)\times (-11)}{(+10)\times (+15)} $ \\ 
    
    	$\square\;$ négatif\qquad $\square\;$ positif\qquad \\
 \end{minipage}
\end{enumerate}
\end{multicols}
\end{EXO}

% @see : https://coopmaths.fr/alea?uuid=2fbc0&id=4C10-8&n=4&d=10&s=10&s2=2&s3=2&s4=true&alea=uCzH&cd=1&cols=2
\begin{EXO}{}{4C10-8}

Trouver les nombres à mettre dans les cases vides pour que les produits de chaque ligne et chaque colonne soient exacts.\\Compléter chaque grille avec des nombres qui conviennent.\\\begin{multicols}{2}
\begin{enumerate}[itemsep=1em]
	\item \begin{minipage}[t]{\linewidth} $\renewcommand{\arraystretch}{2}
\begin{array}{|c|c|c|c|}
\hline
\cellcolor{lightgray} \times & \cellcolor{lightgray} \text{Colonne 1} & \cellcolor{lightgray} \text{Colonne 2} & \cellcolor{lightgray} \text{Produits}\\
\hline
\cellcolor{lightgray} \text{Ligne 1} &  &  & \cellcolor{lightgray}20\\
\hline
 \cellcolor{lightgray} \text{Ligne 2} & -10 &  & \cellcolor{lightgray}80\\
\hline
 \cellcolor{lightgray} \text{Produits} & \cellcolor{lightgray}-40 & \cellcolor{lightgray}-40 & \cellcolor{lightgray}///////\\
\hline
 \end{array}
\renewcommand{\arraystretch}{1}$
 \end{minipage}
	\item \begin{minipage}[t]{\linewidth} $\renewcommand{\arraystretch}{2}
\begin{array}{|c|c|c|c|}
\hline
\cellcolor{lightgray} \times & \cellcolor{lightgray} \text{Colonne 1} & \cellcolor{lightgray} \text{Colonne 2} & \cellcolor{lightgray} \text{Produits}\\
\hline
\cellcolor{lightgray} \text{Ligne 1} &  &  & \cellcolor{lightgray}-16\\
\hline
 \cellcolor{lightgray} \text{Ligne 2} &  & 3 & \cellcolor{lightgray}3\\
\hline
 \cellcolor{lightgray} \text{Produits} & \cellcolor{lightgray}4 & \cellcolor{lightgray}-12 & \cellcolor{lightgray}///////\\
\hline
 \end{array}
\renewcommand{\arraystretch}{1}$
 \end{minipage}
	\item \begin{minipage}[t]{\linewidth} $\renewcommand{\arraystretch}{2}
\begin{array}{|c|c|c|c|}
\hline
\cellcolor{lightgray} \times & \cellcolor{lightgray} \text{Colonne 1} & \cellcolor{lightgray} \text{Colonne 2} & \cellcolor{lightgray} \text{Produits}\\
\hline
\cellcolor{lightgray} \text{Ligne 1} &  & -6 & \cellcolor{lightgray}42\\
\hline
 \cellcolor{lightgray} \text{Ligne 2} &  &  & \cellcolor{lightgray}5\\
\hline
 \cellcolor{lightgray} \text{Produits} & \cellcolor{lightgray}-35 & \cellcolor{lightgray}-6 & \cellcolor{lightgray}///////\\
\hline
 \end{array}
\renewcommand{\arraystretch}{1}$
 \end{minipage}
	\item \begin{minipage}[t]{\linewidth} $\renewcommand{\arraystretch}{2}
\begin{array}{|c|c|c|c|}
\hline
\cellcolor{lightgray} \times & \cellcolor{lightgray} \text{Colonne 1} & \cellcolor{lightgray} \text{Colonne 2} & \cellcolor{lightgray} \text{Produits}\\
\hline
\cellcolor{lightgray} \text{Ligne 1} &  & -10 & \cellcolor{lightgray}-90\\
\hline
 \cellcolor{lightgray} \text{Ligne 2} &  &  & \cellcolor{lightgray}56\\
\hline
 \cellcolor{lightgray} \text{Produits} & \cellcolor{lightgray}72 & \cellcolor{lightgray}-70 & \cellcolor{lightgray}///////\\
\hline
 \end{array}
\renewcommand{\arraystretch}{1}$
 \end{minipage}
\end{enumerate}
\end{multicols}
\end{EXO}

% @see : https://coopmaths.fr/alea?uuid=9e862&id=4C10-9&n=2&d=10&s=10&s2=5&s3=1&alea=OCQS&cd=1&cols=2
\begin{EXO}{}{4C10-9}

Les nombres situés à l'extrémité des flèches sont les produits des nombres dont les flèches sont issues.\begin{multicols}{2}
\begin{enumerate}[itemsep=2em]
	\item \begin{minipage}[t]{\linewidth} Calculer les produits à l'extrémité des flèches.\\\begin{tikzpicture}[baseline,scale = 0.6]

    \tikzset{
      point/.style={
        thick,
        draw,
        cross out,
        inner sep=0pt,
        minimum width=5pt,
        minimum height=5pt,
      },
    }
    \clip (-6.206339097770921,-6.5) rectangle (6.206339097770922,5.354101966249686);
    	
	\draw[color ={black}] (0,0)--(2.28,-0.74);
	\draw[color ={black}] (2.28,-0.74)--(1.76,-2.43);
	\draw[color ={black}] (1.76,-2.43)--(0,-2.4);
	\draw[color ={black}, densely dash dot dot ,-{Stealth[width=5mm]}] (1.76,-2.43)--(3.8,-1.83);
	\draw[color ={black}, densely dash dot dot ,-{Stealth[width=5mm]}] (2.85,0.93)--(4.15,-0.76);
	\draw [color={black}] (1.175571,-1.618034) node[anchor = center,scale=1, rotate = 0] {-4};
	\draw [color={black}] (4.755283,-1.545085) node[anchor = center,scale=1, rotate = 0] {a};
	\draw[color={black}] (5.706339,-1.854102)--(4.755283,-2.545085)--(3.804226,-1.854102)--(4.167497,-0.736068)--(5.343068,-0.736068)--cycle;
	
	\draw[color ={black}] (0,0)--(1.41,1.94);
	\draw[color ={black}] (1.41,1.94)--(2.85,0.93);
	\draw[color ={black}] (2.85,0.93)--(2.28,-0.74);
	\draw[color ={black}, densely dash dot dot ,-{Stealth[width=5mm]}] (2.85,0.93)--(2.91,3.05);
	\draw[color ={black}, densely dash dot dot ,-{Stealth[width=5mm]}] (0,3)--(2,3.71);
	\draw [color={black}] (1.902113,0.618034) node[anchor = center,scale=1, rotate = 0] {3};
	\draw [color={black}] (2.938926,4.045085) node[anchor = center,scale=1, rotate = 0] {b};
	\draw[color={black}] (3.889983,3.736068)--(2.938926,3.045085)--(1.98787,3.736068)--(2.351141,4.854102)--(3.526712,4.854102)--cycle;
	
	\draw[color ={black}] (0,0)--(-1.41,1.94);
	\draw[color ={black}] (-1.41,1.94)--(0,3);
	\draw[color ={black}] (0,3)--(1.41,1.94);
	\draw[color ={black}, densely dash dot dot ,-{Stealth[width=5mm]}] (0,3)--(-2,3.71);
	\draw[color ={black}, densely dash dot dot ,-{Stealth[width=5mm]}] (-2.85,0.93)--(-2.91,3.05);
	\draw [color={black}] (0,2) node[anchor = center,scale=1, rotate = 0] {-8};
	\draw [color={black}] (-2.938926,4.045085) node[anchor = center,scale=1, rotate = 0] {c};
	\draw[color={black}] (-1.98787,3.736068)--(-2.938926,3.045085)--(-3.889983,3.736068)--(-3.526712,4.854102)--(-2.351141,4.854102)--cycle;
	
	\draw[color ={black}] (0,0)--(-2.28,-0.74);
	\draw[color ={black}] (-2.28,-0.74)--(-2.85,0.93);
	\draw[color ={black}] (-2.85,0.93)--(-1.41,1.94);
	\draw[color ={black}, densely dash dot dot ,-{Stealth[width=5mm]}] (-2.85,0.93)--(-4.15,-0.76);
	\draw[color ={black}, densely dash dot dot ,-{Stealth[width=5mm]}] (-1.76,-2.43)--(-3.8,-1.83);
	\draw [color={black}] (-1.902113,0.618034) node[anchor = center,scale=1, rotate = 0] {-6};
	\draw [color={black}] (-4.755283,-1.545085) node[anchor = center,scale=1, rotate = 0] {d};
	\draw[color={black}] (-3.804226,-1.854102)--(-4.755283,-2.545085)--(-5.706339,-1.854102)--(-5.343068,-0.736068)--(-4.167497,-0.736068)--cycle;
	
	\draw[color ={black}] (0,0)--(0,-2.4);
	\draw[color ={black}] (0,-2.4)--(-1.76,-2.43);
	\draw[color ={black}] (-1.76,-2.43)--(-2.28,-0.74);
	\draw[color ={black}, densely dash dot dot ,-{Stealth[width=5mm]}] (-1.76,-2.43)--(-0.57,-4.18);
	\draw[color ={black}, densely dash dot dot ,-{Stealth[width=5mm]}] (1.76,-2.43)--(0.57,-4.18);
	\draw [color={black}] (-1.175571,-1.618034) node[anchor = center,scale=1, rotate = 0] {-3};
	\draw [color={black}] (0,-5) node[anchor = center,scale=1, rotate = 0] {e};
	\draw[color={black}] (0.951057,-5.309017)--(0,-6)--(-0.951057,-5.309017)--(-0.587785,-4.190983)--(0.587785,-4.190983)--cycle;

\end{tikzpicture}\\
 \end{minipage}
	\item \begin{minipage}[t]{\linewidth} Calculer les produits à l'extrémité des flèches.\\\begin{tikzpicture}[baseline,scale = 0.6]

    \tikzset{
      point/.style={
        thick,
        draw,
        cross out,
        inner sep=0pt,
        minimum width=5pt,
        minimum height=5pt,
      },
    }
    \clip (-6.206339097770921,-6.5) rectangle (6.206339097770922,5.354101966249686);
    	
	\draw[color ={black}] (0,0)--(2.28,-0.74);
	\draw[color ={black}] (2.28,-0.74)--(1.76,-2.43);
	\draw[color ={black}] (1.76,-2.43)--(0,-2.4);
	\draw[color ={black}, densely dash dot dot ,-{Stealth[width=5mm]}] (1.76,-2.43)--(3.8,-1.83);
	\draw[color ={black}, densely dash dot dot ,-{Stealth[width=5mm]}] (2.85,0.93)--(4.15,-0.76);
	\draw [color={black}] (1.175571,-1.618034) node[anchor = center,scale=1, rotate = 0] {-3};
	\draw [color={black}] (4.755283,-1.545085) node[anchor = center,scale=1, rotate = 0] {a};
	\draw[color={black}] (5.706339,-1.854102)--(4.755283,-2.545085)--(3.804226,-1.854102)--(4.167497,-0.736068)--(5.343068,-0.736068)--cycle;
	
	\draw[color ={black}] (0,0)--(1.41,1.94);
	\draw[color ={black}] (1.41,1.94)--(2.85,0.93);
	\draw[color ={black}] (2.85,0.93)--(2.28,-0.74);
	\draw[color ={black}, densely dash dot dot ,-{Stealth[width=5mm]}] (2.85,0.93)--(2.91,3.05);
	\draw[color ={black}, densely dash dot dot ,-{Stealth[width=5mm]}] (0,3)--(2,3.71);
	\draw [color={black}] (1.902113,0.618034) node[anchor = center,scale=1, rotate = 0] {5};
	\draw [color={black}] (2.938926,4.045085) node[anchor = center,scale=1, rotate = 0] {b};
	\draw[color={black}] (3.889983,3.736068)--(2.938926,3.045085)--(1.98787,3.736068)--(2.351141,4.854102)--(3.526712,4.854102)--cycle;
	
	\draw[color ={black}] (0,0)--(-1.41,1.94);
	\draw[color ={black}] (-1.41,1.94)--(0,3);
	\draw[color ={black}] (0,3)--(1.41,1.94);
	\draw[color ={black}, densely dash dot dot ,-{Stealth[width=5mm]}] (0,3)--(-2,3.71);
	\draw[color ={black}, densely dash dot dot ,-{Stealth[width=5mm]}] (-2.85,0.93)--(-2.91,3.05);
	\draw [color={black}] (0,2) node[anchor = center,scale=1, rotate = 0] {-9};
	\draw [color={black}] (-2.938926,4.045085) node[anchor = center,scale=1, rotate = 0] {c};
	\draw[color={black}] (-1.98787,3.736068)--(-2.938926,3.045085)--(-3.889983,3.736068)--(-3.526712,4.854102)--(-2.351141,4.854102)--cycle;
	
	\draw[color ={black}] (0,0)--(-2.28,-0.74);
	\draw[color ={black}] (-2.28,-0.74)--(-2.85,0.93);
	\draw[color ={black}] (-2.85,0.93)--(-1.41,1.94);
	\draw[color ={black}, densely dash dot dot ,-{Stealth[width=5mm]}] (-2.85,0.93)--(-4.15,-0.76);
	\draw[color ={black}, densely dash dot dot ,-{Stealth[width=5mm]}] (-1.76,-2.43)--(-3.8,-1.83);
	\draw [color={black}] (-1.902113,0.618034) node[anchor = center,scale=1, rotate = 0] {-6};
	\draw [color={black}] (-4.755283,-1.545085) node[anchor = center,scale=1, rotate = 0] {d};
	\draw[color={black}] (-3.804226,-1.854102)--(-4.755283,-2.545085)--(-5.706339,-1.854102)--(-5.343068,-0.736068)--(-4.167497,-0.736068)--cycle;
	
	\draw[color ={black}] (0,0)--(0,-2.4);
	\draw[color ={black}] (0,-2.4)--(-1.76,-2.43);
	\draw[color ={black}] (-1.76,-2.43)--(-2.28,-0.74);
	\draw[color ={black}, densely dash dot dot ,-{Stealth[width=5mm]}] (-1.76,-2.43)--(-0.57,-4.18);
	\draw[color ={black}, densely dash dot dot ,-{Stealth[width=5mm]}] (1.76,-2.43)--(0.57,-4.18);
	\draw [color={black}] (-1.175571,-1.618034) node[anchor = center,scale=1, rotate = 0] {-2};
	\draw [color={black}] (0,-5) node[anchor = center,scale=1, rotate = 0] {e};
	\draw[color={black}] (0.951057,-5.309017)--(0,-6)--(-0.951057,-5.309017)--(-0.587785,-4.190983)--(0.587785,-4.190983)--cycle;

\end{tikzpicture}\\
 \end{minipage}
\end{enumerate}
\end{multicols}
\end{EXO}

% @see : https://coopmaths.fr/alea?uuid=857c1&id=4C10-10&n=10&d=10&s=10&s2=false&alea=GhuS&cd=1&cols=2
\begin{EXO}{Compléter :}{4C10-10}
\begin{multicols}{2}
\begin{enumerate}[itemsep=2em]
	\item \begin{minipage}[t]{\linewidth} $ \ldots\ldots\ldots \times (-1) = (-3) $ \end{minipage}
	\item \begin{minipage}[t]{\linewidth} $ (-4) \times \ldots\ldots\ldots = (+32) $ \end{minipage}
	\item \begin{minipage}[t]{\linewidth} $ \ldots\ldots\ldots \times (-7) = (+35) $ \end{minipage}
	\item \begin{minipage}[t]{\linewidth} $ \ldots\ldots\ldots \times (+8) = (-32) $ \end{minipage}
	\item \begin{minipage}[t]{\linewidth} $ (+9) \times \ldots\ldots\ldots = (-72) $ \end{minipage}
	\item \begin{minipage}[t]{\linewidth} $ (+9) \times \ldots\ldots\ldots = (-45) $ \end{minipage}
	\item \begin{minipage}[t]{\linewidth} $ \ldots\ldots\ldots \times (-5) = (-25) $ \end{minipage}
	\item \begin{minipage}[t]{\linewidth} $ \ldots\ldots\ldots \times (-10) = (+70) $ \end{minipage}
	\item \begin{minipage}[t]{\linewidth} $ (-6) \times \ldots\ldots\ldots = (-54) $ \end{minipage}
	\item \begin{minipage}[t]{\linewidth} $ \ldots\ldots\ldots \times (-5) = (+30) $ \end{minipage}
\end{enumerate}
\end{multicols}
\end{EXO}


\clearpage

\begin{Correction}
\begin{EXO}{}{}

 $\def\arraystretch{1.5}\begin{array}{|c|c|c|c|c|}
    \hline
    \times & -7 & +3 & -6 & +9 \\
    \hline
    -9 & +63 & -27 & +54 & -81 \\
    \hline
    +5 & -35 & +15 & -30 & +45 \\
    \hline
    -2 & +14 & -6 & +12 & -18 \\
    \hline
    +3 & -21 & +9 & -18 & +27 \\
    \hline
    \end{array}$
\end{EXO}

\begin{EXO}{}{}

 $\def\arraystretch{1.5}\begin{array}{|c|c|c|c|c|}
    \hline
    \times & +7 & -4 & +8 & -3 \\
    \hline
    +3 & +21 & -12 & +24 & -9 \\
    \hline
    -4 & -28 & +16 & -32 & +12 \\
    \hline
    +6 & +42 & -24 & +48 & -18 \\
    \hline
    -2 & -14 & +8 & -16 & +6 \\
    \hline
    \end{array}$
\end{EXO}

\begin{EXO}{}{}

\begin{enumerate}[itemsep=1em]
	\item Il faut que $ y $ soit négatif pour que $A$ soit négatif.
	\item Il faut que $ x $ soit positif pour que $B$ soit positif.
	\item Il faut que $ a $ soit négatif pour que $C$ soit négatif.
	\item Il faut que $ n $ soit négatif pour que $D$ soit positif.
\end{enumerate}
\end{EXO}

\begin{EXO}{}{}

\begin{enumerate}[itemsep=1em]
	\item Il faut que $ n $ soit positif pour que $A$ soit négatif.
	\item Il faut que $ m $ soit négatif pour que $B$ soit positif.
	\item Il faut que $ a $ soit positif pour que $C$ soit positif.
	\item Il faut que $ m $ soit négatif pour que $D$ soit positif.
\end{enumerate}
\end{EXO}

\begin{EXO}{}{}

\begin{enumerate}[itemsep=1em]
	\item Il faut que $ a $ soit positif pour que $A$ soit positif.
	\item Il faut que $ m $ soit négatif pour que $B$ soit positif.
	\item Il faut que $ n $ soit négatif pour que $C$ soit négatif.
	\item Il faut que $ n $ soit positif pour que $D$ soit négatif.
\end{enumerate}
\end{EXO}

\begin{EXO}{}{}

\begin{enumerate}[itemsep=1em]
	\item $\renewcommand{\arraystretch}{2}
\begin{array}{|c|c|c|c|}
\hline
\cellcolor{lightgray} \times & \cellcolor{lightgray} \text{Colonne 1} & \cellcolor{lightgray} \text{Colonne 2} & \cellcolor{lightgray} \text{Produits}\\
\hline
\cellcolor{lightgray} \text{Ligne 1} & \cellcolor[HTML]{F15929} 4 & \cellcolor[HTML]{F15929} 5 & \cellcolor{lightgray}20\\
\hline
 \cellcolor{lightgray} \text{Ligne 2} & \cellcolor[HTML]{F15929} -10 & \cellcolor[HTML]{F15929} -8 & \cellcolor{lightgray}80\\
\hline
 \cellcolor{lightgray} \text{Produits} & \cellcolor{lightgray}-40 & \cellcolor{lightgray}-40 & \cellcolor{lightgray}///////\\
\hline
 \end{array}
\renewcommand{\arraystretch}{1}$

	\item $\renewcommand{\arraystretch}{2}
\begin{array}{|c|c|c|c|}
\hline
\cellcolor{lightgray} \times & \cellcolor{lightgray} \text{Colonne 1} & \cellcolor{lightgray} \text{Colonne 2} & \cellcolor{lightgray} \text{Produits}\\
\hline
\cellcolor{lightgray} \text{Ligne 1} & \cellcolor[HTML]{F15929} 4 & \cellcolor[HTML]{F15929} -4 & \cellcolor{lightgray}-16\\
\hline
 \cellcolor{lightgray} \text{Ligne 2} & \cellcolor[HTML]{F15929} 1 & \cellcolor[HTML]{F15929} 3 & \cellcolor{lightgray}3\\
\hline
 \cellcolor{lightgray} \text{Produits} & \cellcolor{lightgray}4 & \cellcolor{lightgray}-12 & \cellcolor{lightgray}///////\\
\hline
 \end{array}
\renewcommand{\arraystretch}{1}$

	\item $\renewcommand{\arraystretch}{2}
\begin{array}{|c|c|c|c|}
\hline
\cellcolor{lightgray} \times & \cellcolor{lightgray} \text{Colonne 1} & \cellcolor{lightgray} \text{Colonne 2} & \cellcolor{lightgray} \text{Produits}\\
\hline
\cellcolor{lightgray} \text{Ligne 1} & \cellcolor[HTML]{F15929} -7 & \cellcolor[HTML]{F15929} -6 & \cellcolor{lightgray}42\\
\hline
 \cellcolor{lightgray} \text{Ligne 2} & \cellcolor[HTML]{F15929} 5 & \cellcolor[HTML]{F15929} 1 & \cellcolor{lightgray}5\\
\hline
 \cellcolor{lightgray} \text{Produits} & \cellcolor{lightgray}-35 & \cellcolor{lightgray}-6 & \cellcolor{lightgray}///////\\
\hline
 \end{array}
\renewcommand{\arraystretch}{1}$

	\item $\renewcommand{\arraystretch}{2}
\begin{array}{|c|c|c|c|}
\hline
\cellcolor{lightgray} \times & \cellcolor{lightgray} \text{Colonne 1} & \cellcolor{lightgray} \text{Colonne 2} & \cellcolor{lightgray} \text{Produits}\\
\hline
\cellcolor{lightgray} \text{Ligne 1} & \cellcolor[HTML]{F15929} 9 & \cellcolor[HTML]{F15929} -10 & \cellcolor{lightgray}-90\\
\hline
 \cellcolor{lightgray} \text{Ligne 2} & \cellcolor[HTML]{F15929} 8 & \cellcolor[HTML]{F15929} 7 & \cellcolor{lightgray}56\\
\hline
 \cellcolor{lightgray} \text{Produits} & \cellcolor{lightgray}72 & \cellcolor{lightgray}-70 & \cellcolor{lightgray}///////\\
\hline
 \end{array}
\renewcommand{\arraystretch}{1}$

\end{enumerate}
\end{EXO}

\begin{EXO}{}{}

\begin{enumerate}[itemsep=1em]
	\item \begin{tikzpicture}[baseline,scale = 0.6]

    \tikzset{
      point/.style={
        thick,
        draw,
        cross out,
        inner sep=0pt,
        minimum width=5pt,
        minimum height=5pt,
      },
    }
    \clip (-6.206339097770921,-6.5) rectangle (6.206339097770922,5.354101966249686);
    	
	\draw[color ={black}] (0,0)--(2.28,-0.74);
	\draw[color ={black}] (2.28,-0.74)--(1.76,-2.43);
	\draw[color ={black}] (1.76,-2.43)--(0,-2.4);
	\draw[color ={black}, densely dash dot dot ,-{Stealth[width=5mm]}] (1.76,-2.43)--(3.8,-1.83);
	\draw[color ={black}, densely dash dot dot ,-{Stealth[width=5mm]}] (2.85,0.93)--(4.15,-0.76);
	\draw [color={black}] (1.175571,-1.618034) node[anchor = center,scale=1, rotate = 0] {-4};
	\draw [color={black}] (4.755283,-1.545085) node[anchor = center,scale=1, rotate = 0] {-12};
	\draw[color={black}] (5.706339,-1.854102)--(4.755283,-2.545085)--(3.804226,-1.854102)--(4.167497,-0.736068)--(5.343068,-0.736068)--cycle;
	
	\draw[color ={black}] (0,0)--(1.41,1.94);
	\draw[color ={black}] (1.41,1.94)--(2.85,0.93);
	\draw[color ={black}] (2.85,0.93)--(2.28,-0.74);
	\draw[color ={black}, densely dash dot dot ,-{Stealth[width=5mm]}] (2.85,0.93)--(2.91,3.05);
	\draw[color ={black}, densely dash dot dot ,-{Stealth[width=5mm]}] (0,3)--(2,3.71);
	\draw [color={black}] (1.902113,0.618034) node[anchor = center,scale=1, rotate = 0] {3};
	\draw [color={black}] (2.938926,4.045085) node[anchor = center,scale=1, rotate = 0] {-24};
	\draw[color={black}] (3.889983,3.736068)--(2.938926,3.045085)--(1.98787,3.736068)--(2.351141,4.854102)--(3.526712,4.854102)--cycle;
	
	\draw[color ={black}] (0,0)--(-1.41,1.94);
	\draw[color ={black}] (-1.41,1.94)--(0,3);
	\draw[color ={black}] (0,3)--(1.41,1.94);
	\draw[color ={black}, densely dash dot dot ,-{Stealth[width=5mm]}] (0,3)--(-2,3.71);
	\draw[color ={black}, densely dash dot dot ,-{Stealth[width=5mm]}] (-2.85,0.93)--(-2.91,3.05);
	\draw [color={black}] (0,2) node[anchor = center,scale=1, rotate = 0] {-8};
	\draw [color={black}] (-2.938926,4.045085) node[anchor = center,scale=1, rotate = 0] {48};
	\draw[color={black}] (-1.98787,3.736068)--(-2.938926,3.045085)--(-3.889983,3.736068)--(-3.526712,4.854102)--(-2.351141,4.854102)--cycle;
	
	\draw[color ={black}] (0,0)--(-2.28,-0.74);
	\draw[color ={black}] (-2.28,-0.74)--(-2.85,0.93);
	\draw[color ={black}] (-2.85,0.93)--(-1.41,1.94);
	\draw[color ={black}, densely dash dot dot ,-{Stealth[width=5mm]}] (-2.85,0.93)--(-4.15,-0.76);
	\draw[color ={black}, densely dash dot dot ,-{Stealth[width=5mm]}] (-1.76,-2.43)--(-3.8,-1.83);
	\draw [color={black}] (-1.902113,0.618034) node[anchor = center,scale=1, rotate = 0] {-6};
	\draw [color={black}] (-4.755283,-1.545085) node[anchor = center,scale=1, rotate = 0] {18};
	\draw[color={black}] (-3.804226,-1.854102)--(-4.755283,-2.545085)--(-5.706339,-1.854102)--(-5.343068,-0.736068)--(-4.167497,-0.736068)--cycle;
	
	\draw[color ={black}] (0,0)--(0,-2.4);
	\draw[color ={black}] (0,-2.4)--(-1.76,-2.43);
	\draw[color ={black}] (-1.76,-2.43)--(-2.28,-0.74);
	\draw[color ={black}, densely dash dot dot ,-{Stealth[width=5mm]}] (-1.76,-2.43)--(-0.57,-4.18);
	\draw[color ={black}, densely dash dot dot ,-{Stealth[width=5mm]}] (1.76,-2.43)--(0.57,-4.18);
	\draw [color={black}] (-1.175571,-1.618034) node[anchor = center,scale=1, rotate = 0] {-3};
	\draw [color={black}] (0,-5) node[anchor = center,scale=1, rotate = 0] {12};
	\draw[color={black}] (0.951057,-5.309017)--(0,-6)--(-0.951057,-5.309017)--(-0.587785,-4.190983)--(0.587785,-4.190983)--cycle;

\end{tikzpicture}\\

	\item \begin{tikzpicture}[baseline,scale = 0.6]

    \tikzset{
      point/.style={
        thick,
        draw,
        cross out,
        inner sep=0pt,
        minimum width=5pt,
        minimum height=5pt,
      },
    }
    \clip (-6.206339097770921,-6.5) rectangle (6.206339097770922,5.354101966249686);
    	
	\draw[color ={black}] (0,0)--(2.28,-0.74);
	\draw[color ={black}] (2.28,-0.74)--(1.76,-2.43);
	\draw[color ={black}] (1.76,-2.43)--(0,-2.4);
	\draw[color ={black}, densely dash dot dot ,-{Stealth[width=5mm]}] (1.76,-2.43)--(3.8,-1.83);
	\draw[color ={black}, densely dash dot dot ,-{Stealth[width=5mm]}] (2.85,0.93)--(4.15,-0.76);
	\draw [color={black}] (1.175571,-1.618034) node[anchor = center,scale=1, rotate = 0] {-3};
	\draw [color={black}] (4.755283,-1.545085) node[anchor = center,scale=1, rotate = 0] {-15};
	\draw[color={black}] (5.706339,-1.854102)--(4.755283,-2.545085)--(3.804226,-1.854102)--(4.167497,-0.736068)--(5.343068,-0.736068)--cycle;
	
	\draw[color ={black}] (0,0)--(1.41,1.94);
	\draw[color ={black}] (1.41,1.94)--(2.85,0.93);
	\draw[color ={black}] (2.85,0.93)--(2.28,-0.74);
	\draw[color ={black}, densely dash dot dot ,-{Stealth[width=5mm]}] (2.85,0.93)--(2.91,3.05);
	\draw[color ={black}, densely dash dot dot ,-{Stealth[width=5mm]}] (0,3)--(2,3.71);
	\draw [color={black}] (1.902113,0.618034) node[anchor = center,scale=1, rotate = 0] {5};
	\draw [color={black}] (2.938926,4.045085) node[anchor = center,scale=1, rotate = 0] {-45};
	\draw[color={black}] (3.889983,3.736068)--(2.938926,3.045085)--(1.98787,3.736068)--(2.351141,4.854102)--(3.526712,4.854102)--cycle;
	
	\draw[color ={black}] (0,0)--(-1.41,1.94);
	\draw[color ={black}] (-1.41,1.94)--(0,3);
	\draw[color ={black}] (0,3)--(1.41,1.94);
	\draw[color ={black}, densely dash dot dot ,-{Stealth[width=5mm]}] (0,3)--(-2,3.71);
	\draw[color ={black}, densely dash dot dot ,-{Stealth[width=5mm]}] (-2.85,0.93)--(-2.91,3.05);
	\draw [color={black}] (0,2) node[anchor = center,scale=1, rotate = 0] {-9};
	\draw [color={black}] (-2.938926,4.045085) node[anchor = center,scale=1, rotate = 0] {54};
	\draw[color={black}] (-1.98787,3.736068)--(-2.938926,3.045085)--(-3.889983,3.736068)--(-3.526712,4.854102)--(-2.351141,4.854102)--cycle;
	
	\draw[color ={black}] (0,0)--(-2.28,-0.74);
	\draw[color ={black}] (-2.28,-0.74)--(-2.85,0.93);
	\draw[color ={black}] (-2.85,0.93)--(-1.41,1.94);
	\draw[color ={black}, densely dash dot dot ,-{Stealth[width=5mm]}] (-2.85,0.93)--(-4.15,-0.76);
	\draw[color ={black}, densely dash dot dot ,-{Stealth[width=5mm]}] (-1.76,-2.43)--(-3.8,-1.83);
	\draw [color={black}] (-1.902113,0.618034) node[anchor = center,scale=1, rotate = 0] {-6};
	\draw [color={black}] (-4.755283,-1.545085) node[anchor = center,scale=1, rotate = 0] {12};
	\draw[color={black}] (-3.804226,-1.854102)--(-4.755283,-2.545085)--(-5.706339,-1.854102)--(-5.343068,-0.736068)--(-4.167497,-0.736068)--cycle;
	
	\draw[color ={black}] (0,0)--(0,-2.4);
	\draw[color ={black}] (0,-2.4)--(-1.76,-2.43);
	\draw[color ={black}] (-1.76,-2.43)--(-2.28,-0.74);
	\draw[color ={black}, densely dash dot dot ,-{Stealth[width=5mm]}] (-1.76,-2.43)--(-0.57,-4.18);
	\draw[color ={black}, densely dash dot dot ,-{Stealth[width=5mm]}] (1.76,-2.43)--(0.57,-4.18);
	\draw [color={black}] (-1.175571,-1.618034) node[anchor = center,scale=1, rotate = 0] {-2};
	\draw [color={black}] (0,-5) node[anchor = center,scale=1, rotate = 0] {6};
	\draw[color={black}] (0.951057,-5.309017)--(0,-6)--(-0.951057,-5.309017)--(-0.587785,-4.190983)--(0.587785,-4.190983)--cycle;

\end{tikzpicture}\\

\end{enumerate}
\end{EXO}

\begin{EXO}{}{}

\begin{enumerate}[itemsep=1em]
	\item $ {\color{blue}\boldsymbol{(+3)}} \times {\color[HTML]{008002}\boldsymbol{(-1)}} = {\color[HTML]{008002}\boldsymbol{(-3)}} $
	\item $ {\color[HTML]{008002}\boldsymbol{(-4)}} \times {\color[HTML]{008002}\boldsymbol{(-8)}} = {\color{blue}\boldsymbol{(+32)}} $
	\item $ {\color[HTML]{008002}\boldsymbol{(-5)}} \times {\color[HTML]{008002}\boldsymbol{(-7)}} = {\color{blue}\boldsymbol{(+35)}} $
	\item $ {\color[HTML]{008002}\boldsymbol{(-4)}} \times {\color{blue}\boldsymbol{(+8)}} = {\color[HTML]{008002}\boldsymbol{(-32)}} $
	\item $ {\color{blue}\boldsymbol{(+9)}} \times {\color[HTML]{008002}\boldsymbol{(-8)}} = {\color[HTML]{008002}\boldsymbol{(-72)}} $
	\item $ {\color{blue}\boldsymbol{(+9)}} \times {\color[HTML]{008002}\boldsymbol{(-5)}} = {\color[HTML]{008002}\boldsymbol{(-45)}} $
	\item $ {\color{blue}\boldsymbol{(+5)}} \times {\color[HTML]{008002}\boldsymbol{(-5)}} = {\color[HTML]{008002}\boldsymbol{(-25)}} $
	\item $ {\color[HTML]{008002}\boldsymbol{(-7)}} \times {\color[HTML]{008002}\boldsymbol{(-10)}} = {\color{blue}\boldsymbol{(+70)}} $
	\item $ {\color[HTML]{008002}\boldsymbol{(-6)}} \times {\color{blue}\boldsymbol{(+9)}} = {\color[HTML]{008002}\boldsymbol{(-54)}} $
	\item $ {\color[HTML]{008002}\boldsymbol{(-6)}} \times {\color[HTML]{008002}\boldsymbol{(-5)}} = {\color{blue}\boldsymbol{(+30)}} $
\end{enumerate}
\end{EXO}

\clearpage
\end{Correction}
\end{document}

% Local Variables:
% TeX-engine: luatex
% End:
